\documentclass[11pt]{report}
% Neuron document style: Charter + Noto Serif CJK KR, light outlined boxes, grayscale.
% Usage: \documentclass[11pt]{report} % Neuron document style: Charter + Noto Serif CJK KR, light outlined boxes, grayscale.
% Usage: \documentclass[11pt]{report} % Neuron document style: Charter + Noto Serif CJK KR, light outlined boxes, grayscale.
% Usage: \documentclass[11pt]{report} \input{neuron-style}  (compile with xelatex)
\usepackage[a4paper,margin=24mm,headheight=14pt]{geometry}
\usepackage{fontspec}
\setmainfont{Charter}[Path=fonts/,Extension=.otf,UprightFont=*-Regular,ItalicFont=*-Italic,
  BoldFont=*-Bold,BoldItalicFont=*-BoldItalic,Ligatures=TeX]
\setsansfont{lmsans10-regular.otf}[BoldFont=lmsans10-bold.otf,ItalicFont=lmsans10-oblique.otf]
\setmonofont{lmmono10-regular.otf}[ItalicFont=lmmono10-italic.otf]
\newfontfamily\krfont{Nanum Myeongjo}[Scale=0.9]
\XeTeXinterchartokenstate=1
\newXeTeXintercharclass\krclass
\count255="AC00 \loop\ifnum\count255<"D7A4 \XeTeXcharclass\count255=\krclass \advance\count255 by 1 \repeat
\count255="3000 \loop\ifnum\count255<"3040 \XeTeXcharclass\count255=\krclass \advance\count255 by 1 \repeat
\count255="3130 \loop\ifnum\count255<"3190 \XeTeXcharclass\count255=\krclass \advance\count255 by 1 \repeat
\count255="1100 \loop\ifnum\count255<"1200 \XeTeXcharclass\count255=\krclass \advance\count255 by 1 \repeat
\XeTeXinterchartoks 0 \krclass = {\begingroup\krfont}
\XeTeXinterchartoks 4095 \krclass = {\begingroup\krfont}
\XeTeXinterchartoks \krclass 0 = {\endgroup}
\XeTeXinterchartoks \krclass 4095 = {\endgroup}
\XeTeXlinebreaklocale "ko"
\XeTeXlinebreakskip=0pt plus 1pt
\usepackage{amsmath,amssymb,amsthm,mathtools,bm}
\usepackage[italic,defaultmathsizes,noasterisk]{mathastext}
\emergencystretch=2.5em

\usepackage{graphicx,booktabs,tabularx,longtable,array,multirow,enumitem,xcolor}
\usepackage{caption}
\captionsetup{font=small,labelfont=bf}
\usepackage{tikz}
\usetikzlibrary{shapes.geometric,arrows.meta,positioning,fit,calc,backgrounds,decorations.pathreplacing}
\usepackage[most]{tcolorbox}
\usepackage{listings}
\usepackage{fancyhdr}
\usepackage[hidelinks]{hyperref}

\setlist{itemsep=1pt,topsep=3pt,parsep=0pt}
\setlength{\parskip}{3pt}
\renewcommand{\arraystretch}{1.12}

\lstdefinestyle{sql}{language=SQL,basicstyle=\ttfamily\small,keywordstyle=\bfseries,
  commentstyle=\itshape\color{black!60},columns=fullflexible,keepspaces=true,
  frame=l,framerule=0.6pt,rulecolor=\color{black!40},xleftmargin=8pt,framexleftmargin=4pt,
  aboveskip=5pt,belowskip=5pt,showstringspaces=false,breaklines=true,
  morekeywords={OVER,PARTITION,RANK,DENSE_RANK,ROW_NUMBER,NTILE,LAG,LEAD,ROWS,RANGE,PRECEDING,FOLLOWING,UNBOUNDED,CURRENT,ROW,NVL,NVL2,REFERENCES,CASCADE,MINUS,INTERSECT,EXISTS,MATERIALIZED,VIEW,REPLACE,RATIO_TO_REPORT,FIRST_VALUE,LAST_VALUE,INTERVAL,NULLS,LAST,FIRST,CHECK,CONSTRAINT,UNIQUE,DEFERRABLE,COMMIT,BEGIN,ROLLBACK,TRANSACTION,INDEX,DEFAULT,ANY,ALL,SOME,ROUND}}
\lstset{style=sql}

\tcbset{
 nb/.style={enhanced,breakable,colback=white,coltitle=black,fonttitle=\bfseries,
   left=3mm,right=3mm,top=1.6mm,bottom=1.6mm,before skip=9pt,after skip=9pt,parbox=false},
 framed/.style={nb,boxrule=0.5pt,colframe=black!70,arc=1.2mm,colbacktitle=white,
   titlerule=0.4pt,titlerule style={black!25},toptitle=1mm,bottomtitle=0.6mm},
 light/.style={framed,colframe=black!35},
 ruled/.style={nb,boxrule=0pt,leftrule=1.1pt,colframe=black!45,sharp corners,left=3.5mm,
   top=0.5mm,bottom=0.5mm,detach title,before upper={\tcbtitle\par\smallskip}},
}
\DeclareTColorBox{keybox}{o}{framed,leftrule=2.2pt,title={\IfValueTF{#1}{#1}{Key idea}}}
\DeclareTColorBox{formulabox}{o}{framed,IfValueTF={#1}{title={#1}}{}}
\DeclareTColorBox{workout}{o}{light,title={Example\IfValueT{#1}{\ \textmd{---}\ #1}}}
\DeclareTColorBox{qa}{m}{light,title={#1}}
\DeclareTColorBox{warnbox}{o}{ruled,leftrule=2.2pt,colframe=black,title={Caution\IfValueT{#1}{\ \textmd{---}\ #1}}}
\DeclareTColorBox{whybox}{o}{ruled,title={\IfValueTF{#1}{#1}{Why it works}}}
\DeclareTColorBox{prof}{o}{ruled,title={Lecture note\IfValueT{#1}{\ \textmd{---}\ #1}}}
\DeclareTColorBox{career}{o}{ruled,title={Career note\IfValueT{#1}{\ \textmd{---}\ #1}}}
\DeclareTColorBox{insight}{o}{ruled,title={In practice\IfValueT{#1}{\ \textmd{---}\ #1}}}

\newcommand{\Step}[1]{\par\smallskip\noindent\textbf{Step #1.}\ }
\newcommand{\Concept}[1]{\hfill{\small[개념: #1]}}
\newcommand{\Use}[1]{\noindent\textit{사용 개념:} #1\par}
\newcommand{\Ans}[1]{\par\smallskip\noindent\fbox{\parbox{\dimexpr\linewidth-2\fboxsep-2\fboxrule\relax}{\textbf{답}\ \ #1}}\par}
\newcommand{\Trap}[1]{\par\smallskip\noindent\textit{흔한 실수:} #1}
\newcommand{\Verify}[1]{\par\smallskip\noindent\textit{검산:} #1}
\newcommand{\Intuit}[1]{\par\smallskip\textbf{직관: #1}}
\newcommand{\pk}[1]{\underline{#1}}
\newcommand{\NULL}{\textsc{null}}
\newcommand{\sqlk}[1]{\texttt{#1}}
\newcommand{\Nul}{{\color{black!50}\textit{null}}}

% relational algebra symbols
\newcommand{\sel}{\sigma}
\newcommand{\proj}{\pi}
\newcommand{\join}{\bowtie}

\pagestyle{fancy}
\fancyhf{}
\fancyhead[R]{\small\itshape\leftmark}
\fancyfoot[C]{\small\thepage}
\renewcommand{\headrulewidth}{0pt}
\renewcommand{\chaptermark}[1]{\markboth{#1}{}}

\usepackage{titlesec}
\titleformat{\chapter}[display]{\bfseries\Large}{\chaptertitlename\ \thechapter}{6pt}{\LARGE}
\titlespacing*{\chapter}{0pt}{-10pt}{14pt}
\titleformat{\section}{\bfseries\large}{\thesection}{0.8em}{}
\titleformat{\subsection}{\bfseries\normalsize}{\thesubsection}{0.8em}{}

% ---------- cheat-sheet / problem-sheet helpers ----------
\newcounter{qn}
\newcommand{\Q}[2]{\par\smallskip\refstepcounter{qn}\noindent\textbf{Q\theqn.}\ {\footnotesize[#1]}\ #2\par}
\newcommand{\A}[1]{\vspace{1pt}\noindent\begin{tcolorbox}[ruled,leftrule=1.6pt,colframe=black!45,before skip=1pt,after skip=4pt,top=0.5mm,bottom=0.5mm]\small\textbf{답} #1\end{tcolorbox}}
\newcommand{\Lab}[3]{\item \textbf{#1}\\ \textit{해 본 것:} #2\\ \textit{\ensuremath{\Rightarrow} 개념:} #3}
\newcommand{\eg}{\textit{예)}\ }
\newcommand{\bi}{\textit{비유)}\ }
  (compile with xelatex)
\usepackage[a4paper,margin=24mm,headheight=14pt]{geometry}
\usepackage{fontspec}
\setmainfont{Charter}[Path=fonts/,Extension=.otf,UprightFont=*-Regular,ItalicFont=*-Italic,
  BoldFont=*-Bold,BoldItalicFont=*-BoldItalic,Ligatures=TeX]
\setsansfont{lmsans10-regular.otf}[BoldFont=lmsans10-bold.otf,ItalicFont=lmsans10-oblique.otf]
\setmonofont{lmmono10-regular.otf}[ItalicFont=lmmono10-italic.otf]
\newfontfamily\krfont{Nanum Myeongjo}[Scale=0.9]
\XeTeXinterchartokenstate=1
\newXeTeXintercharclass\krclass
\count255="AC00 \loop\ifnum\count255<"D7A4 \XeTeXcharclass\count255=\krclass \advance\count255 by 1 \repeat
\count255="3000 \loop\ifnum\count255<"3040 \XeTeXcharclass\count255=\krclass \advance\count255 by 1 \repeat
\count255="3130 \loop\ifnum\count255<"3190 \XeTeXcharclass\count255=\krclass \advance\count255 by 1 \repeat
\count255="1100 \loop\ifnum\count255<"1200 \XeTeXcharclass\count255=\krclass \advance\count255 by 1 \repeat
\XeTeXinterchartoks 0 \krclass = {\begingroup\krfont}
\XeTeXinterchartoks 4095 \krclass = {\begingroup\krfont}
\XeTeXinterchartoks \krclass 0 = {\endgroup}
\XeTeXinterchartoks \krclass 4095 = {\endgroup}
\XeTeXlinebreaklocale "ko"
\XeTeXlinebreakskip=0pt plus 1pt
\usepackage{amsmath,amssymb,amsthm,mathtools,bm}
\usepackage[italic,defaultmathsizes,noasterisk]{mathastext}
\emergencystretch=2.5em

\usepackage{graphicx,booktabs,tabularx,longtable,array,multirow,enumitem,xcolor}
\usepackage{caption}
\captionsetup{font=small,labelfont=bf}
\usepackage{tikz}
\usetikzlibrary{shapes.geometric,arrows.meta,positioning,fit,calc,backgrounds,decorations.pathreplacing}
\usepackage[most]{tcolorbox}
\usepackage{listings}
\usepackage{fancyhdr}
\usepackage[hidelinks]{hyperref}

\setlist{itemsep=1pt,topsep=3pt,parsep=0pt}
\setlength{\parskip}{3pt}
\renewcommand{\arraystretch}{1.12}

\lstdefinestyle{sql}{language=SQL,basicstyle=\ttfamily\small,keywordstyle=\bfseries,
  commentstyle=\itshape\color{black!60},columns=fullflexible,keepspaces=true,
  frame=l,framerule=0.6pt,rulecolor=\color{black!40},xleftmargin=8pt,framexleftmargin=4pt,
  aboveskip=5pt,belowskip=5pt,showstringspaces=false,breaklines=true,
  morekeywords={OVER,PARTITION,RANK,DENSE_RANK,ROW_NUMBER,NTILE,LAG,LEAD,ROWS,RANGE,PRECEDING,FOLLOWING,UNBOUNDED,CURRENT,ROW,NVL,NVL2,REFERENCES,CASCADE,MINUS,INTERSECT,EXISTS,MATERIALIZED,VIEW,REPLACE,RATIO_TO_REPORT,FIRST_VALUE,LAST_VALUE,INTERVAL,NULLS,LAST,FIRST,CHECK,CONSTRAINT,UNIQUE,DEFERRABLE,COMMIT,BEGIN,ROLLBACK,TRANSACTION,INDEX,DEFAULT,ANY,ALL,SOME,ROUND}}
\lstset{style=sql}

\tcbset{
 nb/.style={enhanced,breakable,colback=white,coltitle=black,fonttitle=\bfseries,
   left=3mm,right=3mm,top=1.6mm,bottom=1.6mm,before skip=9pt,after skip=9pt,parbox=false},
 framed/.style={nb,boxrule=0.5pt,colframe=black!70,arc=1.2mm,colbacktitle=white,
   titlerule=0.4pt,titlerule style={black!25},toptitle=1mm,bottomtitle=0.6mm},
 light/.style={framed,colframe=black!35},
 ruled/.style={nb,boxrule=0pt,leftrule=1.1pt,colframe=black!45,sharp corners,left=3.5mm,
   top=0.5mm,bottom=0.5mm,detach title,before upper={\tcbtitle\par\smallskip}},
}
\DeclareTColorBox{keybox}{o}{framed,leftrule=2.2pt,title={\IfValueTF{#1}{#1}{Key idea}}}
\DeclareTColorBox{formulabox}{o}{framed,IfValueTF={#1}{title={#1}}{}}
\DeclareTColorBox{workout}{o}{light,title={Example\IfValueT{#1}{\ \textmd{---}\ #1}}}
\DeclareTColorBox{qa}{m}{light,title={#1}}
\DeclareTColorBox{warnbox}{o}{ruled,leftrule=2.2pt,colframe=black,title={Caution\IfValueT{#1}{\ \textmd{---}\ #1}}}
\DeclareTColorBox{whybox}{o}{ruled,title={\IfValueTF{#1}{#1}{Why it works}}}
\DeclareTColorBox{prof}{o}{ruled,title={Lecture note\IfValueT{#1}{\ \textmd{---}\ #1}}}
\DeclareTColorBox{career}{o}{ruled,title={Career note\IfValueT{#1}{\ \textmd{---}\ #1}}}
\DeclareTColorBox{insight}{o}{ruled,title={In practice\IfValueT{#1}{\ \textmd{---}\ #1}}}

\newcommand{\Step}[1]{\par\smallskip\noindent\textbf{Step #1.}\ }
\newcommand{\Concept}[1]{\hfill{\small[개념: #1]}}
\newcommand{\Use}[1]{\noindent\textit{사용 개념:} #1\par}
\newcommand{\Ans}[1]{\par\smallskip\noindent\fbox{\parbox{\dimexpr\linewidth-2\fboxsep-2\fboxrule\relax}{\textbf{답}\ \ #1}}\par}
\newcommand{\Trap}[1]{\par\smallskip\noindent\textit{흔한 실수:} #1}
\newcommand{\Verify}[1]{\par\smallskip\noindent\textit{검산:} #1}
\newcommand{\Intuit}[1]{\par\smallskip\textbf{직관: #1}}
\newcommand{\pk}[1]{\underline{#1}}
\newcommand{\NULL}{\textsc{null}}
\newcommand{\sqlk}[1]{\texttt{#1}}
\newcommand{\Nul}{{\color{black!50}\textit{null}}}

% relational algebra symbols
\newcommand{\sel}{\sigma}
\newcommand{\proj}{\pi}
\newcommand{\join}{\bowtie}

\pagestyle{fancy}
\fancyhf{}
\fancyhead[R]{\small\itshape\leftmark}
\fancyfoot[C]{\small\thepage}
\renewcommand{\headrulewidth}{0pt}
\renewcommand{\chaptermark}[1]{\markboth{#1}{}}

\usepackage{titlesec}
\titleformat{\chapter}[display]{\bfseries\Large}{\chaptertitlename\ \thechapter}{6pt}{\LARGE}
\titlespacing*{\chapter}{0pt}{-10pt}{14pt}
\titleformat{\section}{\bfseries\large}{\thesection}{0.8em}{}
\titleformat{\subsection}{\bfseries\normalsize}{\thesubsection}{0.8em}{}

% ---------- cheat-sheet / problem-sheet helpers ----------
\newcounter{qn}
\newcommand{\Q}[2]{\par\smallskip\refstepcounter{qn}\noindent\textbf{Q\theqn.}\ {\footnotesize[#1]}\ #2\par}
\newcommand{\A}[1]{\vspace{1pt}\noindent\begin{tcolorbox}[ruled,leftrule=1.6pt,colframe=black!45,before skip=1pt,after skip=4pt,top=0.5mm,bottom=0.5mm]\small\textbf{답} #1\end{tcolorbox}}
\newcommand{\Lab}[3]{\item \textbf{#1}\\ \textit{해 본 것:} #2\\ \textit{\ensuremath{\Rightarrow} 개념:} #3}
\newcommand{\eg}{\textit{예)}\ }
\newcommand{\bi}{\textit{비유)}\ }
  (compile with xelatex)
\usepackage[a4paper,margin=24mm,headheight=14pt]{geometry}
\usepackage{fontspec}
\setmainfont{Charter}[Path=fonts/,Extension=.otf,UprightFont=*-Regular,ItalicFont=*-Italic,
  BoldFont=*-Bold,BoldItalicFont=*-BoldItalic,Ligatures=TeX]
\setsansfont{lmsans10-regular.otf}[BoldFont=lmsans10-bold.otf,ItalicFont=lmsans10-oblique.otf]
\setmonofont{lmmono10-regular.otf}[ItalicFont=lmmono10-italic.otf]
\newfontfamily\krfont{Nanum Myeongjo}[Scale=0.9]
\XeTeXinterchartokenstate=1
\newXeTeXintercharclass\krclass
\count255="AC00 \loop\ifnum\count255<"D7A4 \XeTeXcharclass\count255=\krclass \advance\count255 by 1 \repeat
\count255="3000 \loop\ifnum\count255<"3040 \XeTeXcharclass\count255=\krclass \advance\count255 by 1 \repeat
\count255="3130 \loop\ifnum\count255<"3190 \XeTeXcharclass\count255=\krclass \advance\count255 by 1 \repeat
\count255="1100 \loop\ifnum\count255<"1200 \XeTeXcharclass\count255=\krclass \advance\count255 by 1 \repeat
\XeTeXinterchartoks 0 \krclass = {\begingroup\krfont}
\XeTeXinterchartoks 4095 \krclass = {\begingroup\krfont}
\XeTeXinterchartoks \krclass 0 = {\endgroup}
\XeTeXinterchartoks \krclass 4095 = {\endgroup}
\XeTeXlinebreaklocale "ko"
\XeTeXlinebreakskip=0pt plus 1pt
\usepackage{amsmath,amssymb,amsthm,mathtools,bm}
\usepackage[italic,defaultmathsizes,noasterisk]{mathastext}
\emergencystretch=2.5em

\usepackage{graphicx,booktabs,tabularx,longtable,array,multirow,enumitem,xcolor}
\usepackage{caption}
\captionsetup{font=small,labelfont=bf}
\usepackage{tikz}
\usetikzlibrary{shapes.geometric,arrows.meta,positioning,fit,calc,backgrounds,decorations.pathreplacing}
\usepackage[most]{tcolorbox}
\usepackage{listings}
\usepackage{fancyhdr}
\usepackage[hidelinks]{hyperref}

\setlist{itemsep=1pt,topsep=3pt,parsep=0pt}
\setlength{\parskip}{3pt}
\renewcommand{\arraystretch}{1.12}

\lstdefinestyle{sql}{language=SQL,basicstyle=\ttfamily\small,keywordstyle=\bfseries,
  commentstyle=\itshape\color{black!60},columns=fullflexible,keepspaces=true,
  frame=l,framerule=0.6pt,rulecolor=\color{black!40},xleftmargin=8pt,framexleftmargin=4pt,
  aboveskip=5pt,belowskip=5pt,showstringspaces=false,breaklines=true,
  morekeywords={OVER,PARTITION,RANK,DENSE_RANK,ROW_NUMBER,NTILE,LAG,LEAD,ROWS,RANGE,PRECEDING,FOLLOWING,UNBOUNDED,CURRENT,ROW,NVL,NVL2,REFERENCES,CASCADE,MINUS,INTERSECT,EXISTS,MATERIALIZED,VIEW,REPLACE,RATIO_TO_REPORT,FIRST_VALUE,LAST_VALUE,INTERVAL,NULLS,LAST,FIRST,CHECK,CONSTRAINT,UNIQUE,DEFERRABLE,COMMIT,BEGIN,ROLLBACK,TRANSACTION,INDEX,DEFAULT,ANY,ALL,SOME,ROUND}}
\lstset{style=sql}

\tcbset{
 nb/.style={enhanced,breakable,colback=white,coltitle=black,fonttitle=\bfseries,
   left=3mm,right=3mm,top=1.6mm,bottom=1.6mm,before skip=9pt,after skip=9pt,parbox=false},
 framed/.style={nb,boxrule=0.5pt,colframe=black!70,arc=1.2mm,colbacktitle=white,
   titlerule=0.4pt,titlerule style={black!25},toptitle=1mm,bottomtitle=0.6mm},
 light/.style={framed,colframe=black!35},
 ruled/.style={nb,boxrule=0pt,leftrule=1.1pt,colframe=black!45,sharp corners,left=3.5mm,
   top=0.5mm,bottom=0.5mm,detach title,before upper={\tcbtitle\par\smallskip}},
}
\DeclareTColorBox{keybox}{o}{framed,leftrule=2.2pt,title={\IfValueTF{#1}{#1}{Key idea}}}
\DeclareTColorBox{formulabox}{o}{framed,IfValueTF={#1}{title={#1}}{}}
\DeclareTColorBox{workout}{o}{light,title={Example\IfValueT{#1}{\ \textmd{---}\ #1}}}
\DeclareTColorBox{qa}{m}{light,title={#1}}
\DeclareTColorBox{warnbox}{o}{ruled,leftrule=2.2pt,colframe=black,title={Caution\IfValueT{#1}{\ \textmd{---}\ #1}}}
\DeclareTColorBox{whybox}{o}{ruled,title={\IfValueTF{#1}{#1}{Why it works}}}
\DeclareTColorBox{prof}{o}{ruled,title={Lecture note\IfValueT{#1}{\ \textmd{---}\ #1}}}
\DeclareTColorBox{career}{o}{ruled,title={Career note\IfValueT{#1}{\ \textmd{---}\ #1}}}
\DeclareTColorBox{insight}{o}{ruled,title={In practice\IfValueT{#1}{\ \textmd{---}\ #1}}}

\newcommand{\Step}[1]{\par\smallskip\noindent\textbf{Step #1.}\ }
\newcommand{\Concept}[1]{\hfill{\small[개념: #1]}}
\newcommand{\Use}[1]{\noindent\textit{사용 개념:} #1\par}
\newcommand{\Ans}[1]{\par\smallskip\noindent\fbox{\parbox{\dimexpr\linewidth-2\fboxsep-2\fboxrule\relax}{\textbf{답}\ \ #1}}\par}
\newcommand{\Trap}[1]{\par\smallskip\noindent\textit{흔한 실수:} #1}
\newcommand{\Verify}[1]{\par\smallskip\noindent\textit{검산:} #1}
\newcommand{\Intuit}[1]{\par\smallskip\textbf{직관: #1}}
\newcommand{\pk}[1]{\underline{#1}}
\newcommand{\NULL}{\textsc{null}}
\newcommand{\sqlk}[1]{\texttt{#1}}
\newcommand{\Nul}{{\color{black!50}\textit{null}}}

% relational algebra symbols
\newcommand{\sel}{\sigma}
\newcommand{\proj}{\pi}
\newcommand{\join}{\bowtie}

\pagestyle{fancy}
\fancyhf{}
\fancyhead[R]{\small\itshape\leftmark}
\fancyfoot[C]{\small\thepage}
\renewcommand{\headrulewidth}{0pt}
\renewcommand{\chaptermark}[1]{\markboth{#1}{}}

\usepackage{titlesec}
\titleformat{\chapter}[display]{\bfseries\Large}{\chaptertitlename\ \thechapter}{6pt}{\LARGE}
\titlespacing*{\chapter}{0pt}{-10pt}{14pt}
\titleformat{\section}{\bfseries\large}{\thesection}{0.8em}{}
\titleformat{\subsection}{\bfseries\normalsize}{\thesubsection}{0.8em}{}

% ---------- cheat-sheet / problem-sheet helpers ----------
\newcounter{qn}
\newcommand{\Q}[2]{\par\smallskip\refstepcounter{qn}\noindent\textbf{Q\theqn.}\ {\footnotesize[#1]}\ #2\par}
\newcommand{\A}[1]{\vspace{1pt}\noindent\begin{tcolorbox}[ruled,leftrule=1.6pt,colframe=black!45,before skip=1pt,after skip=4pt,top=0.5mm,bottom=0.5mm]\small\textbf{답} #1\end{tcolorbox}}
\newcommand{\Lab}[3]{\item \textbf{#1}\\ \textit{해 본 것:} #2\\ \textit{\ensuremath{\Rightarrow} 개념:} #3}
\newcommand{\eg}{\textit{예)}\ }
\newcommand{\bi}{\textit{비유)}\ }

\usepackage{float}
\newcommand{\ft}[1]{\texttt{\detokenize{#1}}}
\newcolumntype{L}{>{\raggedright\arraybackslash}X}
\lstdefinestyle{py}{language=Python,basicstyle=\ttfamily\small,keywordstyle=\bfseries,commentstyle=\itshape\color{black!55},
  frame=leftline,framerule=0.8pt,rulecolor=\color{black!45},xleftmargin=3mm,showstringspaces=false,breaklines=true,columns=fullflexible,
  numbers=left,numberstyle=\tiny\color{black!50},numbersep=4pt}
\lstdefinestyle{out}{basicstyle=\ttfamily\footnotesize,frame=single,rulecolor=\color{black!25},xleftmargin=3mm,columns=fullflexible,breaklines=true,numbers=none}
\lstset{style=py}
\begin{document}

\begin{titlepage}\centering
\vspace*{30mm}
{\Huge\bfseries 코드 따라 읽기\par}\medskip
{\Large LendingClub 승인 전략의 핵심 코드를 한 줄씩\par}
\vspace{12mm}
{\large 서울대 ABS 14기 팀 프로젝트 내부용\par}
\vspace{4mm}
{\small 2026년 11월 16일 \ensuremath{\cdot} 코드는 저장소의 \ft{src/lc/}와 \ft{scripts/}에서 그대로 가져왔다(영어 주석, 한글 이름표는 영어로 바꿈)\par}
\vfill
{\small 처음이면 1--2장부터, 파이썬을 조금 알면 3장부터 읽는다. 각 장의 ``따라 해 보기''는 \ft{docs/notes/code/snippets.py}로 바로 실행된다.\par}
\end{titlepage}

\tableofcontents
\chapter{시작 전에: 지도와 실행법}

\section{이 문서를 읽는 법}
코드를 처음부터 끝까지 외울 필요는 없다. 이 프로젝트의 코드는 결국 \textbf{여섯 단계를 한 줄로 이은 것}이다. 각 장은 한 단계를 맡고, 같은 순서로 쓴다.
\begin{enumerate}[leftmargin=*]
\item \textbf{무엇을 하려는가}(한 문단).
\item \textbf{실제 코드}(줄 번호 포함).
\item \textbf{줄별 설명}(``3번 줄'' 식으로 짚는다).
\item \textbf{따라 해 보기}: 데이터 없이 도는 작은 예제와 그 출력.
\item \textbf{흔한 실수}.
\end{enumerate}

\begin{keybox}[전체 흐름 한 장]
\small\begin{tabularx}{\linewidth}{@{}p{5mm}Lp{40mm}@{}}\toprule
단계 & 하는 일 & 파일\\\midrule
1 & CSV를 읽어 상태를 ``상환/부도''로 정리, 미래 정보 제거 & \ft{src/lc/data.py}\\
2 & 그 시점에 결과를 알 수 있던 대출만 학습에 쓰기 & \ft{src/lc/asof.py}\\
3 & LightGBM으로 부도확률 $p$ 만들기 & scripts 23, 30번\\
4 & $p$를 기대수익 점수로 바꾸고 발행월마다 상위 $q$ 승인 & scripts 30, 38, 39번\\
5 & 승인한 묶음의 실제 현금흐름으로 IRR 채점 & \ft{src/lc/returns.py}\\
6 & 매달 초과수익 \ensuremath{\to} Sharpe, 부트스트랩, 1-SE 선택, 파레토 & scripts 39, 41번\\\bottomrule
\end{tabularx}
\Intuit{예상(점수)으로 고르고, 실제 결과(IRR)로 채점하고, 그 채점표를 48개월 모아 Sharpe를 낸다.}
\end{keybox}

\section{폴더 구조}
\begin{table}[H]\centering\small
\begin{tabularx}{\linewidth}{@{}p{38mm}L@{}}\toprule
위치 & 들어 있는 것\\\midrule
\ft{src/lc/} & 여러 스크립트가 같이 쓰는 \textbf{함수 모음}(라이브러리). 데이터 정리, 시점 학습, 수익률, 비교 기준\\
\ft{scripts/NN_이름.py} & 실험 하나 = 파일 하나. 번호가 실험 순서다. 실행하면 \ft{results/}에 표와 로그를 남긴다\\
\ft{results/} & 실행 결과(\ft{.csv}, \ft{.log}). 문서의 숫자는 모두 여기서 나온다\\
\ft{docs/} & 결정 기록(\ft{decisions.md}), 연구 노트(\ft{research_log.md}), test 규칙(\ft{test_protocol*.md}), 보고서\\
\ft{data/} & 원본 CSV, 변환한 parquet, 국채·\allowbreak{}회사채 금리(\ft{external/})\\\bottomrule
\end{tabularx}
\end{table}

\section{실행하는 법}
터미널(맥은 ``터미널'' 앱)을 열고 저장소 폴더로 이동한 뒤 실행한다. 이 컴퓨터는 conda 환경 \ft{cmapss}의 파이썬을 쓴다.
\begin{lstlisting}[language=bash,numbers=none]
cd ~/Desktop/learningClub
conda activate cmapss                       # once per terminal window
python scripts/39_final_choice.py           # one experiment (a few minutes)
python docs/notes/code/snippets.py irr      # one practice example (instant)
\end{lstlisting}
\begin{warnbox}[python을 못 찾는다고 나오면]
\ft{command not found: python}은 conda 환경을 켜지 않았다는 뜻이다. \ft{conda activate cmapss}를 먼저 하거나, 전체 경로 \ft{/opt/homebrew/Caskroom/miniconda/base/envs/cmapss/bin/python}으로 실행한다. 맥 기본 \ft{python3}에는 LightGBM이 없다.
\end{warnbox}

\section{스크립트 맨 위의 두 가지 특이한 줄}
거의 모든 스크립트가 이렇게 시작한다.
\begin{lstlisting}
import _path  # noqa: F401
from pathlib import Path

SRC = (Path(__file__).parent / '30_portfolio_sharpe.py').read_text()
exec(SRC[:SRC.index("sharpe = lambda")])
\end{lstlisting}
\begin{itemize}[leftmargin=*]
\item \textbf{1번 줄} \ft{import _path}: \ft{scripts/_path.py}는 ``\ft{src/} 폴더도 찾아봐라''라고 파이썬에 알려 주는 세 줄짜리 파일이다. 그래서 다음 줄부터 \ft{from lc import returns}처럼 라이브러리를 불러올 수 있다. \ft{# noqa}는 ``안 쓰는 import라고 경고하지 마''라는 표시일 뿐이다.
\item \textbf{4번 줄}: 30번 스크립트의 \textbf{내용을 글자로} 읽는다. \ft{Path(__file__).parent}는 ``지금 이 파일이 있는 폴더''다.
\item \textbf{5번 줄}: 그 글자 중 \ft{"sharpe = lambda"}가 나오기 \textbf{직전까지만} 잘라(\ft{SRC[:위치]}) 실행한다(\ft{exec}). 30번의 앞부분은 데이터를 읽고 모델을 학습해 \ft{ev} 표를 만드는 부분이다. 결국 ``30번이 만든 \ft{ev}를 그대로 이어받아 쓴다''는 뜻이다.
\end{itemize}
\begin{warnbox}[왜 이렇게 하나, 무엇을 조심하나]
같은 학습 코드를 복사하지 않으려는 지름길이다. 대신 \textbf{앞 스크립트의 변수 이름이 갑자기 나타난다}(\ft{ev}, \ft{months}, \ft{TB} 등). 처음 보는 변수가 있으면 \ft{exec}한 파일에서 찾는다. 39번은 30번을, 41번은 39번을, 40번은 35번을, 35번은 15번을 이어받는다. 또 앞 파일에서 자르는 기준 문자열이 바뀌면 \ft{ValueError: substring not found}가 난다.
\end{warnbox}

\chapter{파이썬, numpy, pandas: 이 코드에 나오는 것만}

이 장은 문법 교과서가 아니다. 프로젝트 코드에 \textbf{실제로 나오는 것만} 모았다. 모르는 기호가 나오면 이 장으로 돌아온다.

\section{파이썬 기본 여섯 가지}
\begin{lstlisting}
QS = [0.1, 0.2, 0.3]                  # list: ordered values, QS[0] is 0.1
hurdles = {'tb': 0.015, 'baa': 0.05}  # dict: name -> value, hurdles['tb']
for q in QS:                          # loop: q takes each value in turn
    print(f'approve {q:.0%}')         # f-string: {q:.0%} prints 10%
sharpe = lambda x: x.mean() / x.std(ddof=1)   # one-line function
def cut(score, q):                    # normal function
    return score >= 0.5
\end{lstlisting}
\begin{itemize}[leftmargin=*]
\item \textbf{1번 줄} 리스트: 대괄호, 순서가 있다. 번호는 0부터 센다.
\item \textbf{2번 줄} 딕셔너리: 중괄호, ``이름: 값'' 짝이다. \ft{hurdles.items()}는 (이름, 값) 짝을 하나씩 준다.
\item \textbf{3--4번 줄} \ft{for}: 들여쓰기(공백 4칸)한 줄들이 반복된다. 파이썬은 \textbf{들여쓰기가 문법}이라, 칸이 틀리면 오류가 난다.
\item \textbf{4번 줄} f-문자열: 따옴표 앞 \ft{f}, 중괄호 안에 변수. \ft{:.2f}는 소수 둘째 자리, \texttt{:.0\%}는 백분율이다.
\item \textbf{5번 줄} \ft{lambda}: 이름 없는 한 줄 함수. ``\ft{x}를 받아 평균 \ensuremath{\div} 표준편차를 돌려준다''.
\item \textbf{6--7번 줄} \ft{def}: 여러 줄 함수. \ft{return} 뒤가 결과다.
\end{itemize}
그 밖에 자주 나오는 것:
\begin{itemize}[leftmargin=*]
\item \ft{[f(x) for x in xs]}: ``xs의 각 x에 f를 적용한 리스트''(리스트 컴프리헨션).
\item \ft{a, b = 1, 2}: 한 번에 두 변수.
\item \ft{x if cond else y}: 조건에 따라 x 또는 y.
\item \ft{**dict}: 딕셔너리를 펼쳐 넣기. \texttt{\{'a': 1, **row\}}는 \ft{row}의 내용을 합친 새 딕셔너리다.
\end{itemize}

\section{numpy: 숫자 배열을 한 번에}
\ft{np.array}는 숫자를 줄지어 담은 배열이다. 핵심은 \textbf{반복문 없이 전체에 한 번에 계산한다}는 것이다.
\begin{lstlisting}
import numpy as np
r = np.array([0.07, 0.11, 0.15])
bad = np.array([0, 1, 0])
r * 100                      # [7, 11, 15]  every element at once
np.where(bad == 1, 0.0, r)   # [0.07, 0, 0.15]  if default then 0 else rate
np.quantile(r, 0.7)          # 70th percentile
r >= np.quantile(r, 0.7)     # [False, False, True]  boolean array
r[r > 0.1]                   # [0.11, 0.15]  keep where True
rng = np.random.default_rng(11)   # random generator with a fixed seed
rng.random(3)                # 3 numbers in [0, 1), same every run
\end{lstlisting}
\begin{itemize}[leftmargin=*]
\item \textbf{4번 줄}: 배열 \ensuremath{\times} 숫자는 원소마다 곱한다.
\item \textbf{5번 줄} \ft{np.where(조건, 참일 때, 거짓일 때)}: 엑셀의 IF를 배열 전체에 한 번에 한다.
\item \textbf{6--7번 줄} 분위수: \ft{np.quantile(s, 1 - q)} 이상이면 ``상위 $q$''다. 상위 30\%를 고르려면 70\% 분위수 이상을 고른다. 이 프로젝트의 모든 승인 규칙이 이 한 줄이다.
\item \textbf{8번 줄}: 참/거짓 배열을 대괄호에 넣으면 참인 자리만 남는다.
\item \textbf{9--10번 줄}: seed(씨앗)를 고정하면 ``무작위''도 매번 같은 값이 나와 결과를 재현할 수 있다. 코드의 \ft{default_rng(2026)}, \ft{default_rng(500 + k)}가 모두 이것이다.
\end{itemize}
\paragraph{자주 나오는 함수.} \ft{x.mean()} 평균, \ft{x.std(ddof=1)} 표본 표준편차(\ft{ddof=1}은 $n-1$로 나눈다는 뜻), \ft{x.sum()}, \ft{np.zeros(n, bool)} 거짓 $n$개, \ft{np.ones}, \ft{np.arange(5)}는 0--4, \ft{np.isclose(a, b)}는 소수 비교(0.8과 0.80000001을 같다고 본다).
\paragraph{2차원 배열과 \ft{axis}.} \ft{x[BIX]}처럼 번호 표로 뽑으면 표 모양 배열이 된다. \ft{.mean(axis=1)}은 \textbf{행마다} 평균, \ft{axis=0}은 \textbf{열마다} 평균이다. 10장의 부트스트랩에서 쓴다.

\section{pandas: 엑셀 표를 코드로}
\ft{pd.DataFrame}은 행과 열이 있는 표, 열 하나는 \ft{Series}다.
\begin{lstlisting}
import pandas as pd
df['r']                          # one column (Series)
df[['id', 'r']]                  # several columns (DataFrame)
df['bad'] = (df['status'] == 'Charged Off').astype(int)   # new column
df[df['grade'] == 'A']           # rows where condition is True
df.loc[rows, 'r']                # by label     df.iloc[3]  by position
df['r'].to_numpy()               # Series -> numpy array
df.groupby('ym')['amt'].sum()    # sum per issue month
df.groupby('ym').indices         # {month: array of row positions}
d.merge(tb, on=['ym', 'term_m'], how='left')   # VLOOKUP
\end{lstlisting}
\begin{itemize}[leftmargin=*]
\item \textbf{4번 줄}: 비교(\ft{==})는 참/거짓 열을 만들고, \ft{.astype(int)}가 1/0으로 바꾼다. 부도 여부 \ft{bad}가 이렇게 만들어진다.
\item \textbf{5번 줄}: 조건을 여러 개 묶을 때는 \texttt{\&}(그리고), \ft{|}(또는), \texttt{\textasciitilde}(아님)을 쓰고 \textbf{각 조건을 괄호로 감싼다}. \ft{and}/\ft{or}를 쓰면 오류가 난다.
\item \textbf{6번 줄} \ft{loc}은 이름표(인덱스) 기준, \ft{iloc}은 0부터 센 위치 기준이다. 코드가 \ft{reset_index(drop=True)}를 자주 하는 이유는 이름표를 0, 1, 2, \ldots{}로 다시 매겨 둘을 같게 만들기 위해서다.
\item \textbf{7번 줄}: 계산을 빠르게 하려고 열을 numpy 배열로 바꾼다. 예를 들어 \ft{P_amt}는 \ft{funded_amnt} 열을 배열로 바꾼 것이다.
\item \textbf{9번 줄} \ft{groupby().indices}: \textbf{이 프로젝트에서 가장 중요한 pandas 표현}이다. ``발행월마다 그 달 대출의 행 번호''를 준다. 발행월마다 자르기(8장), 발행월마다 IRR(9장)이 모두 이것을 돈다.
\item \textbf{10번 줄} \ft{merge}: 엑셀의 VLOOKUP. 발행월과 만기가 같은 국채 금리를 대출 옆에 붙인다(6장).
\end{itemize}

\begin{workout}[따라 해 보기: \ft{snippets.py groupby}]
6건짜리 표로 달마다 합계, 등급별 부도율, 달마다 행 번호를 본다. 출력의 마지막 두 줄:
\begin{lstlisting}[style=out]
2013-01 -> rows [0 1 2]
2013-02 -> rows [3 4 5]
\end{lstlisting}
\Ans{1월 대출은 0--2번 행, 2월 대출은 3--5번 행이다. 이후 모든 ``발행월마다''는 이 행 번호 묶음을 하나씩 도는 것이다.}
\end{workout}
\Trap{\ft{df[df['r'] > 0.1 and df['grade'] == 'A']}. pandas에서는 \texttt{(df['r'] > 0.1) \& (df['grade'] == 'A')}로 써야 한다.}

\chapter{데이터 준비와 시점 학습}

\section{상태 정리: \texttt{data.py}}
\textbf{하려는 것.} 원본의 \ft{loan_status}는 ``Fully Paid'', ``Charged Off'', ``Late (31-120 days)'' 등 문자열이다. 이것을 상환(paid), 부도(default), 연체(late), 진행 중(current) 네 가지로 정리하고, 숫자로 쓸 열(금리, 만기, 발행월)을 만든다.
\begin{lstlisting}
def classify(status):
    s = status.astype(str)
    return pd.Series(np.select(
        [s.str.contains('Fully Paid'), s.str.contains('Charged Off') | s.eq('Default'),
         s.str.startswith('Late') | s.eq('In Grace Period')],
        ['paid', 'default', 'late'], 'current'), index=status.index)

def parse(df):
    d = df.copy()
    d['ym'] = pd.to_datetime(d['issue_d'], format='%b-%Y')
    d['issue_year'] = d['ym'].dt.year
    d['r'] = pd.to_numeric(d['int_rate'].astype('string').str.strip().str.rstrip('%')) / 100
    d['term_m'] = d['term'].astype('string').str.extract(r'(\d+)')[0].astype(int)
    lp = pd.to_datetime(d['last_pymnt_d'], format='%b-%Y', errors='coerce')
    d['tau'] = months_between(d['ym'], lp).fillna(1).clip(lower=1).astype(int)
    d['outcome'] = classify(d['loan_status'])
    d['resolved'] = d['outcome'].isin(['paid', 'default'])
    d['bad'] = (d['outcome'] == 'default').astype('int8')
    return d
\end{lstlisting}
\begin{itemize}[leftmargin=*]
\item \textbf{3--6번 줄} \ft{np.select([조건1, 조건2, 조건3], [값1, 값2, 값3], 기본값)}: 엑셀의 중첩 IF. 위에서부터 처음 참인 조건의 값을 쓰고, 아무것도 아니면 \ft{'current'}다. \ft{.str.contains}는 ``그 글자를 포함하나'', \ft{.str.startswith}는 ``그 글자로 시작하나''다.
\item \textbf{9번 줄} \ft{df.copy()}: 원본을 건드리지 않으려고 복사본에서 작업한다.
\item \textbf{10번 줄}: ``Dec-2015'' 같은 글자를 날짜로 바꾼다. \texttt{\%b}는 월 약자, \texttt{\%Y}는 네 자리 연도다. 이 \ft{ym}(발행월)이 이후 모든 ``발행월마다''의 기준이다.
\item \textbf{12번 줄}: ``13.56\%'' \ensuremath{\to} 공백 제거 \ensuremath{\to} \% 제거 \ensuremath{\to} 숫자 13.56 \ensuremath{\to} 100으로 나눠 0.1356.
\item \textbf{13번 줄}: `` 36 months''에서 숫자만 뽑는다. \texttt{r'(\textbackslash d+)'}는 ``연속된 숫자''라는 정규식이다.
\item \textbf{14--15번 줄} \ft{tau}: 발행부터 마지막 납부까지 개월 수. 현금흐름 복원(7장)과 시점 학습(아래)에 쓴다. \ft{errors='coerce'}는 ``변환이 안 되면 빈 값으로''이고, \ft{fillna(1)}은 빈 값을 1로 채운다.
\item \textbf{17--18번 줄}: 결과가 확정된(상환 또는 부도) 대출만 \ft{resolved}=참. 부도면 \ft{bad}=1이 모델의 정답이다.
\end{itemize}

\paragraph{누수 목록 \ft{LEAK}.} \ft{total_pymnt}(받은 총액), \ft{recoveries}(회수금), \ft{last_pymnt_d}처럼 \textbf{승인한 뒤에야 알 수 있는} 열들이다. \ft{feature_columns()}가 이것을 모델 입력에서 뺀다. 이 열들은 채점(IRR)에는 쓰지만 고르기(모델)에는 절대 쓰지 않는다.
\begin{warnbox}[누수]
\ft{total_pymnt}를 모델에 넣으면 AUC가 0.99 가까이 나온다. 이미 결과를 보고 맞히는 것이라 실전에서는 쓸 수 없다. 새 열을 넣을 때마다 ``승인하는 날 이 값을 알 수 있나?''를 묻는다.
\end{warnbox}

\section{시점 학습: \texttt{asof.py}}
\textbf{하려는 것.} 2016년 1월 1일에 승인 결정을 한다면, 그날까지 \textbf{결과가 확정된} 대출로만 배워야 한다. 2015년 9월에 나간 대출이 2016년 3월에 부도났다는 사실은 2016년 1월에는 모른다.
\begin{lstlisting}
CHARGE_OFF_LAG = 5

def months_to(d, cutoff):
    c = pd.Timestamp(cutoff)
    return ((c.year - d['ym'].dt.year) * 12 + c.month - d['ym'].dt.month).to_numpy()

def known_at(d, cutoff, lag=CHARGE_OFF_LAG):
    age = months_to(d, cutoff)
    settle = d['tau'].to_numpy() + np.where(d['bad'].to_numpy() == 1, lag, 0)
    return (age > 0) & d['resolved'].to_numpy() & (settle <= age)

def training_set(d, cutoff, lag=CHARGE_OFF_LAG):
    return d[(d['ym'] < pd.Timestamp(cutoff)).to_numpy() & known_at(d, cutoff, lag)]
\end{lstlisting}
\begin{itemize}[leftmargin=*]
\item \textbf{1번 줄}: 부도는 마지막 납부 뒤 약 150일(5개월) 연체 후에 상각(Charged Off)되므로, 5개월 늦게 확정된다고 본다.
\item \textbf{5번 줄} \ft{age}: 발행월부터 기준일까지 개월 수. (연도 차이 \ensuremath{\times} 12) + 월 차이.
\item \textbf{9번 줄} \ft{settle}: 결과가 확정된 개월. 상환이면 \ft{tau}, 부도면 \ft{tau} + 5.
\item \textbf{10번 줄}: 세 조건을 모두 만족해야 학습에 쓸 수 있다. 기준일 전에 발행됐고, 결과가 확정됐고, 확정 시점이 기준일 이전이다.
\end{itemize}
\begin{workout}[따라 해 보기: \ft{snippets.py asof}]
기준일 2016-01-01, 대출 4건:
\begin{lstlisting}[style=out]
          ym  tau  bad  resolved  age  settle  known
0 2014-01-01   12    0      True   24      12   True
1 2014-06-01    6    1      True   19      11   True
2 2015-03-01    3    0      True   10       3   True
3 2015-09-01    2    1      True    4       7  False
\end{lstlisting}
\Step{1} 3번 대출: 2015년 9월 발행, 2개월 납부 뒤 부도. 확정은 $2+5=7$개월째인데 기준일까지는 4개월뿐이다.
\Ans{3번은 2016년 1월에는 ``아직 모르는'' 대출이라 학습에서 빠진다. 이 규칙이 없으면 미래를 보고 배우는 셈이다.}
\end{workout}
\Trap{기준일 이전에 \emph{발행된} 대출을 모두 쓰는 것. 그러면 빨리 끝난 대출(대부분 조기상환)만 남거나, 아직 모를 부도가 섞인다. 둘 다 결과를 왜곡한다(보고서 3장의 종결 편향).}

\chapter{부도확률 모델과 기대수익 점수}

\section{LightGBM 학습: \texttt{fit\_lgb}}
\textbf{하려는 것.} 신청 정보(특성 \ft{X})로 부도 여부(\ft{y}=\ft{bad})를 맞히는 나무 모델을 학습하고, 평가 연도 대출의 부도확률 $p$를 낸다.
\begin{lstlisting}
LGB_PARAMS = dict(n_estimators=2000, learning_rate=0.05, num_leaves=63, subsample=0.8, subsample_freq=1,
                  colsample_bytree=0.6, min_child_samples=200, reg_lambda=5, cat_smooth=50, n_jobs=-1, verbose=-1)

def fit_lgb(Xt, yt, seed, extra=None):
    p = dict(models.LGB_PARAMS)
    p.update(extra or {})
    u = np.random.default_rng(seed).random(len(Xt)) < 0.15
    m = lgb.LGBMClassifier(random_state=seed, **p)
    m.fit(Xt[~u], yt[~u], eval_set=[(Xt[u], yt[u])], callbacks=[lgb.early_stopping(50, verbose=False)])
    return m
\end{lstlisting}
\begin{itemize}[leftmargin=*]
\item \textbf{1--2번 줄} 설정값(고정, 튜닝 안 함). 핵심만 보면 이렇다.
  \begin{itemize}
  \item \ft{n_estimators=2000}: 나무를 최대 2000그루까지 쌓는다.
  \item \ft{learning_rate=0.05}: 나무 하나가 앞 나무들의 실수를 5\%씩만 고친다. 작을수록 천천히, 안정적으로 배운다.
  \item \ft{num_leaves=63}: 나무 하나가 손님을 최대 63개 무리로 나눈다.
  \item \ft{min_child_samples=200}: 한 무리에 최소 200건. 몇 건에 맞춘 규칙(과적합)을 막는다.
  \item \ft{subsample}, \ft{colsample_bytree}: 나무마다 행 80\%, 열 60\%만 무작위로 써서 나무들이 서로 다르게 배우게 한다.
  \end{itemize}
\item \textbf{7번 줄}: 학습 표본의 15\%를 무작위로 떼어(\ft{u}=참) 검증용으로 둔다.
\item \textbf{9번 줄}: 나머지 85\%(\texttt{\textasciitilde u})로 학습한다. 검증 15\%의 성적이 50그루 연속 나아지지 않으면 멈춘다(조기 종료). 2000은 상한일 뿐이고 실제로는 수백 그루에서 멈춘다.
\end{itemize}

\section{연도마다 학습하고 $p$, 점수 만들기: 30번의 앞부분}
\begin{lstlisting}
EV = []
for Y in range(2013, 2017):
    tr = asof.training_set(d, f'{Y}-01-01')
    tri, ytr = tr.index, tr['bad'].to_numpy()
    usable = [c for c in feats if X0.loc[tri, c].nunique(dropna=True) > 1]
    evi = d.index[(d['issue_year'] == Y) & d['resolved'] & (d['term_m'] == 36)]
    ev = returns.attach_treasury(d.loc[evi]).reset_index(drop=True)
    lgd = asof.loss_given_default(tr)
    p = np.mean([fit_lgb(X0.loc[tri, usable], ytr, s).predict_proba(X0.loc[evi, usable])[:, 1]
                 for s in SEEDS], axis=0)
    ev['p'] = p
    ev['s_er'] = (1 - p) * ev['r'] - p * lgd - ev['tb']
    EV.append(ev)
ev = pd.concat(EV, ignore_index=True)
\end{lstlisting}
\begin{itemize}[leftmargin=*]
\item \textbf{2번 줄}: 평가 연도 2013, 2014, 2015, 2016을 하나씩 돈다(\ft{range(2013, 2017)}은 2017을 포함하지 않는다).
\item \textbf{3번 줄}: 그해 1월 1일 기준으로 알 수 있던 대출만으로 학습 표본을 만든다(2장).
\item \textbf{5번 줄}: 학습 표본에서 값이 한 가지뿐인 열(예: 아직 생기지 않은 컬럼)은 뺀다.
\item \textbf{6번 줄} 평가 대상: 그해 발행된, 결과가 확정된, 36개월 대출.
\item \textbf{7번 줄}: 국채 금리 \ft{tb}를 붙이고 행 번호를 0부터 다시 매긴다.
\item \textbf{8번 줄} LGD: 학습 표본 부도 대출의 평균 손실률, 즉 $1-$(받은 총액/대출액)의 평균이다.
\item \textbf{9--10번 줄}: seed 0, 1, 2로 모델을 세 번 학습해 $p$를 평균 낸다. 무작위 요소(검증 분할, 행·\allowbreak{}열 추출) 때문에 생기는 흔들림을 줄인다. \ft{predict_proba}는 [부도 아닐 확률, 부도일 확률] 두 열을 주고, \ft{[:, 1]}이 두 번째 열, 즉 부도확률이다.
\item \textbf{12번 줄} \textbf{기대수익 점수}: $s=(1-p)\,r-p\cdot\mathrm{LGD}-tb$. 안 망하면 받는 이자, 빼기 망하면 잃는 원금, 빼기 국채 기회비용.
\item \textbf{14번 줄}: 4개 연도 표를 위아래로 붙여 하나의 \ft{ev}로 만든다. 이후 모든 장의 \ft{ev}가 이것이다.
\end{itemize}

\begin{workout}[따라 해 보기: \ft{snippets.py lgbm}]
가짜 데이터 5000건(부도율 약 17\%)으로 같은 설정을 seed 세 개로 학습한다.
\begin{lstlisting}[style=out]
seed 0: 158 trees, AUC 0.982
seed 1: 182 trees, AUC 0.985
seed 2: 76 trees, AUC 0.985
3-seed mean AUC 0.984, mean p 0.161 vs actual 0.166
\end{lstlisting}
\Ans{2000그루 상한이지만 조기 종료로 76--182그루에서 멈췄다. seed마다 나무 수가 다르고, 평균을 내면 조금 더 안정적이다. 가짜 데이터라 AUC가 높고, 실제 LC 데이터에서는 0.63--0.68이다.}
\end{workout}
\Trap{\ft{predict_proba(X)}를 그대로 $p$로 쓰는 것. 두 열짜리 표이므로 \ft{[:, 1]}을 빠뜨리면 모양이 맞지 않는다. 또 최신 LightGBM은 \ft{eval_set} 대신 \ft{eval_X}, \ft{eval_y}를 권한다는 경고를 낼 수 있는데, 결과는 같다.}

\section{국채 붙이기: \texttt{attach\_treasury}}
\begin{lstlisting}
def attach_treasury(d):
    m = d.merge(treasury(), on=['ym', 'term_m'], how='left', validate='many_to_one')
    m.index = d.index
    assert m['tb'].notna().all()
    return m
\end{lstlisting}
\ft{treasury()}는 FRED에서 받은 GS3(3년), GS5(5년) 월별 금리 표다. \textbf{2번 줄}은 발행월과 만기가 같은 금리를 VLOOKUP한다. \ft{validate='many_to_one'}은 ``대출 여러 건에 금리 한 줄''이 맞는지 검사한다. \textbf{4번 줄} \ft{assert}는 ``빈 금리가 하나라도 있으면 즉시 멈춰라''라는 안전장치다.

\chapter{채점: 현금흐름, IRR, 고르기}

\section{현금흐름 복원: \texttt{schedule}, \texttt{cashflows}}
\textbf{하려는 것.} 데이터에는 월별 납부 기록이 없고 ``받은 총액''(\ft{total_pymnt}), ``월 상환액''(\ft{installment}), ``마지막 납부월''(\ft{tau})만 있다. 이것으로 ``몇 달은 정상 납부, 마지막 달에 나머지 일시금''이라고 근사한다.
\begin{lstlisting}
def schedule(d):
    inst = d['installment'].to_numpy(float)
    tot = d['total_pymnt'].to_numpy(float)
    n = d['tau'].to_numpy(int)
    kf = np.minimum(n, np.floor(tot / np.maximum(inst, 1e-9))).astype(int)
    rem = np.maximum(tot - kf * inst, 0.0)
    return kf, rem, n
\end{lstlisting}
\begin{itemize}[leftmargin=*]
\item \textbf{5번 줄} \ft{kf}(정상 납부 개월): 받은 총액 \ensuremath{\div} 월 상환액을 내림한 값. 단 마지막 납부월 \ft{n}을 넘을 수 없다. \ft{np.maximum(inst, 1e-9)}는 0으로 나누기를 막는다.
\item \textbf{6번 줄} \ft{rem}(일시금): 받은 총액에서 정상 납부분을 뺀 나머지. 조기상환이면 남은 원금이고, 부도면 회수금이다.
\end{itemize}
예: 월 33.21달러짜리 대출에서 총 500달러를 받고 마지막 납부가 20개월째라면 다음과 같다.
\[ kf=\min(20,\ \lfloor 500/33.21\rfloor)=15,\qquad rem=500-15\times33.21=1.85 \]
즉 15개월은 33.21달러씩, 20개월째에 1.85달러를 받는다.

\ft{cashflows()}는 이것을 \textbf{묶음 전체의 월별 현금흐름 배열}(0--99개월)로 만든다. 0개월째는 대출액 합계의 음수(투자)이고, 1--\ft{kf}개월은 월 상환액, \ft{n}개월째에 일시금이 들어온다. 수수료 1\%는 받은 돈마다 \ft{(1 - fee)}를 곱한다. 코드의 \ft{np.cumsum}, \ft{np.add.at}은 ``대출마다 반복문을 돌지 않고 한 번에 더하기'' 위한 기술이다. 뜻은 위 문장과 같다.

\section{IRR: \texttt{irr}}
\begin{lstlisting}
def irr(cf):
    cf = np.where(np.abs(cf) < 1e-6, 0.0, cf)
    f = lambda m: (cf / (1 + m) ** np.arange(len(cf))).sum()
    try:
        m = brentq(f, -0.2, 0.2)
    except ValueError:
        m = brentq(f, -0.95, 1.0)
    return (1 + m) ** 12 - 1
\end{lstlisting}
\begin{itemize}[leftmargin=*]
\item \textbf{3번 줄} \ft{f(m)}: 월 할인율 $m$으로 모든 현금흐름을 현재가치로 바꿔 더한 값(NPV). $t$개월째 돈은 $(1+m)^t$로 나눈다.
\item \textbf{5번 줄} \ft{brentq(f, a, b)}: $a$와 $b$ 사이에서 $f(m)=0$이 되는 $m$을 찾는다(엑셀 목표값 찾기). NPV를 0으로 만드는 할인율이 IRR이다.
\item \textbf{6--7번 줄}: 손실이 아주 크면 $-20\%$--$20\%$ 범위에 답이 없어서, 범위를 넓혀 다시 찾는다.
\item \textbf{8번 줄}: 월 수익률을 연 수익률로 바꾼다(복리).
\end{itemize}
\begin{workout}[따라 해 보기: \ft{snippets.py irr}]
1000달러, 연 12\%, 36개월.
\begin{lstlisting}[style=out]
monthly payment 33.21
fully paid: IRR 0.1268
1% fee: IRR 0.1191
default after 10: IRR -0.6814
pool of two: IRR -0.1405
\end{lstlisting}
\Step{1} 다 갚으면 12.68\%다. 12\%가 아닌 이유는 월 1\%를 12번 복리로 쌓으면 $1.01^{12}-1=12.68\%$이기 때문이다.
\Step{2} 수수료 1\%가 연 수익률을 약 0.8\%p 깎는다.
\Step{3} 한 건짜리 부도는 $-68\%$처럼 극단값이 된다. 그래서 대출 하나하나의 IRR로 Sharpe를 내지 않는다(보고서 6장).
\Step{4} 두 대출을 \textbf{현금흐름을 합쳐} 한 묶음으로 보면 $-14\%$다. 두 IRR의 평균($(12.68-68.14)/2=-27.7\%$)과 다르다.
\Ans{묶음 IRR = 현금흐름을 먼저 합치고 IRR 한 번. 이 프로젝트의 채점은 모두 ``발행월 묶음 IRR''이다(\ft{pooled_irr}).}
\end{workout}

\section{고르기: 발행월마다 상위 $q$}
\begin{lstlisting}
def cut(score, q, keys):
    a = np.zeros(len(ev), bool)
    for _, ix in ev.groupby(keys, observed=True).indices.items():
        s = score[ix]
        a[ix] = s >= np.quantile(s, 1 - q)
    return a

jit = np.random.default_rng(11).random(len(ev)) * 1e-6
er = ev['s_er'].to_numpy() + jit            # V1: expected return
lowpd = -ev['p'].to_numpy() + jit           # V2: low default probability
a = cut(er, 0.3, 'ym')                      # approve top 30% within each month
a3 = cut(-acc + jit, 0.3, ['ym', 'grade'])  # V3: within month x grade
\end{lstlisting}
\begin{itemize}[leftmargin=*]
\item \textbf{2번 줄}: 승인 여부 배열을 전부 ``거절''로 시작한다.
\item \textbf{3번 줄}: 발행월(\ft{keys='ym'})마다 그 달 대출의 행 번호 \ft{ix}를 돈다. \ft{_}는 ``이름(그 달 날짜)은 안 쓴다''는 표시다.
\item \textbf{4--5번 줄}: 그 달 점수 중 상위 $q$를 승인으로 바꾼다. \textbf{발행월마다} 자르므로 매달 승인 금액 비중이 거의 일정하다. 연도 단위로 자르면 이 비중이 들쭉날쭉해 Sharpe가 왜곡된다(37단계의 교훈).
\item \textbf{8번 줄} 동점 깨기: 점수에 0--0.000001 사이의 아주 작은 무작위 값을 더한다. 순위는 거의 그대로이고, 동점만 무작위로 갈린다.
\item \textbf{10번 줄}: ``부도확률이 낮을수록 좋다''를 ``점수가 높을수록 좋다''로 맞추려고 마이너스를 붙인다.
\item \textbf{12번 줄}: \ft{keys}에 두 열을 주면 ``발행월 \ensuremath{\times} 등급''마다 자른다. 같은 등급 안에서 고르는 규칙이다.
\end{itemize}
\begin{workout}[따라 해 보기: \ft{snippets.py cut}]
\begin{lstlisting}[style=out]
top 50%: [1 1 1 0 0 1 1 0]
ties as is: [1 1 1 1 1 1 1 1] -> approves all, not 50%
ties + tiny noise: [0 1 1 0 1 1 0 0]
\end{lstlisting}
\Step{1} 첫 줄: 1월 점수 0.9, 0.5, 0.5, 0.1에서 50\% 분위수는 0.5다. 0.5 이상이 셋이라 75\%가 승인됐다. 동점이 있으면 정확히 $q$가 안 된다.
\Step{2} 둘째 줄: 전부 동점이면 전부 승인된다. 34단계에서 실제로 이 버그가 나서, 계좌 개설 수(정수) 규칙의 실제 승인비율이 30\%가 아니라 39.5\%였다.
\Ans{셋째 줄처럼 작은 무작위 값을 더하면 매달 정확히 절반씩 승인된다.}
\end{workout}
\Trap{점수가 정수이거나 같은 값이 많으면 꼭 동점을 깬다. 그리고 고른 뒤에는 \ft{a.mean()}으로 실제 승인비율을 확인한다.}

\paragraph{무작위 기준선.} \ft{cut(np.random.default_rng(500 + k).random(len(ev)), q, 'ym')}. 점수 대신 무작위 수로 같은 비율을 고른다. seed 500--509로 10번(학습 쪽은 5번) 해서 평균 낸다. 이것이 ``아무거나 고를 때''다.

\chapter{Sharpe, 불확실성, 최종 선택, test}

\section{발행월 포트폴리오 초과수익: \texttt{port\_series}}
\begin{lstlisting}
def port_series(a, hurdle):
    out = []
    for m in months:
        ix = grp[m]
        sel = ix[a[ix]]
        r = returns.pooled_irr(ev.iloc[sel], fee=FEE) - hurdle[ix].mean() if len(sel) else 0.0
        out.append(P_amt[sel].sum() / P_amt[ix].sum() * r)
    return np.array(out)

sharpe = lambda x: x.mean(axis=-1) / x.std(axis=-1, ddof=1)
\end{lstlisting}
\begin{itemize}[leftmargin=*]
\item \textbf{3--4번 줄}: 48개 발행월을 차례로 돈다. \ft{grp = ev.groupby('ym').indices}(30번에서 만듦)라서 \ft{grp[m]}은 그 달 행 번호다.
\item \textbf{5번 줄}: 그 달 대출 중 승인된 것만 남긴다. \ft{a[ix]}는 그 달 대출들의 승인 여부이고, 그것으로 \ft{ix}를 거른다.
\item \textbf{6번 줄}: 승인 묶음의 IRR(수수료 1\%)에서 그 달 허들(국채 등)을 뺀다. 승인이 하나도 없으면 0이다.
\item \textbf{7번 줄} \textbf{핵심}: 그 달 포트폴리오 초과수익 = (승인 금액 \ensuremath{\div} 그 달 전체 금액) \ensuremath{\times} 초과수익. 거절한 돈은 국채에 넣었다고 보므로 그 부분의 초과수익은 0이다.
\item \textbf{10번 줄}: Sharpe = 48개 값의 평균 \ensuremath{\div} 표준편차. \ft{axis=-1}은 ``마지막 축으로''라는 뜻이다. 1차원이면 그냥 전체이고, 2차원(부트스트랩 표)이면 행마다 계산한다.
\end{itemize}
\begin{workout}[따라 해 보기: \ft{snippets.py sharpe}]
\begin{lstlisting}[style=out]
share 0.3: mean 0.0124, std 0.0053, Sharpe 2.33
share 0.8: mean 0.0331, std 0.0142, Sharpe 2.33
\end{lstlisting}
\Ans{비중을 0.3에서 0.8로 바꾸면 평균과 표준편차가 모두 같은 배수로 커져 Sharpe는 그대로다. 그래서 승인비율 사이의 Sharpe 차이는 ``얼마나 샀나''가 아니라 ``얼마나 까다롭게 골랐나''의 차이다. 같은 이유로 파레토 비교는 같은 승인비율 안에서 한다(아래).}
\end{workout}

\section{이동 블록 부트스트랩: 표준오차와 구간}
\textbf{하려는 것.} 48개월은 짧다. ``다른 48개월이 왔다면 Sharpe가 얼마나 달랐을까''를 알려고, 있는 달들을 다시 뽑아 가상의 48개월을 2000번 만든다. 이웃한 달끼리 비슷하므로(경기 흐름) 한 달씩이 아니라 \textbf{6개월 덩어리}로 뽑는다.
\begin{lstlisting}
n = len(months)                                   # 48
st = np.random.default_rng(2026).integers(0, n - BLOCK + 1, size=(B, int(np.ceil(n / BLOCK))))
BIX = np.concatenate([np.arange(BLOCK)[None, :] + st[:, [j]] for j in range(st.shape[1])], axis=1)[:, :n]

se = sharpe(port[BIX]).std(ddof=1)                # standard error of Sharpe
diff_bs = sharpe(ch[BIX]) - np.mean([sharpe(r[BIX]) for r in rnd], axis=0)
lo, hi = np.quantile(diff_bs, 0.05), np.quantile(diff_bs, 0.95)   # 90% interval
\end{lstlisting}
\begin{itemize}[leftmargin=*]
\item \textbf{2번 줄}: 덩어리 시작점을 무작위로 뽑는다. 크기는 (2000번) \ensuremath{\times} (덩어리 8개)이고, 시작점은 0--42 사이다(덩어리가 48을 넘지 않게).
\item \textbf{3번 줄}: 시작점 $s$를 $s, s+1, \ldots, s+5$로 늘여 이어 붙인다. 결과 \ft{BIX}는 2000행 \ensuremath{\times} 48열의 ``달 번호 표''다. 행 하나가 가상의 48개월 하나다.
\item \textbf{5번 줄}: \ft{port[BIX]}는 2000 \ensuremath{\times} 48 표다. 행마다 Sharpe를 내면 2000개의 Sharpe가 나오고, 그 흩어짐(표준편차)이 \textbf{표준오차(SE)}다.
\item \textbf{6번 줄} \textbf{짝지은 비교}: 같은 가상 48개월(같은 \ft{BIX} 행)에서 선택안과 무작위를 함께 계산해 차이를 낸다. 시장 전체가 좋았던 달 같은 공통 흔들림이 차이에서 지워져, 차이의 구간이 각각의 SE보다 훨씬 좁다.
\item \textbf{7번 줄}: 2000개 차이의 하위 5\%와 상위 5\%가 90\% 구간이다. 하한이 0보다 크면 ``거의 확실히 무작위보다 낫다''.
\end{itemize}
\begin{workout}[따라 해 보기: \ft{snippets.py boot}]
\ft{n=12, block=3, B=3}으로 줄여 본 \ft{BIX}:
\begin{lstlisting}[style=out]
[[ 8  9 10  1  2  3  0  1  2  6  7  8]
 [ 3  4  5  4  5  6  0  1  2  3  4  5]
 [ 6  7  8  3  4  5  8  9 10  7  8  9]]
\end{lstlisting}
\Ans{행마다 세 달씩 연속인 덩어리 네 개가 이어져 있다. 같은 달이 두 번 뽑히기도 하고(복원 추출), 안 뽑히는 달도 있다.}
\end{workout}

\section{1-SE 규칙으로 최종 제안 고르기: 39번}
\begin{lstlisting}
best = C.loc[C['Sharpe'].idxmax()]
ok = C[C['Sharpe'] >= best['Sharpe'] - best['SE']]
pick = ok.sort_values(['q', 'Sharpe'], ascending=[False, False]).iloc[0]
\end{lstlisting}
\begin{itemize}[leftmargin=*]
\item \ft{C}는 후보 18개(전략 2 \ensuremath{\times} 승인비율 9)의 Sharpe와 SE 표다.
\item \textbf{1번 줄} \ft{idxmax()}: Sharpe가 가장 큰 행의 이름표. \ft{loc}으로 그 행을 꺼낸다.
\item \textbf{2번 줄}: (최고 Sharpe $-$ 그 SE) 이상인 후보만 남긴다. 최고와 통계적으로 구분되지 않는 후보들이다.
\item \textbf{3번 줄}: 승인비율 큰 순, 같으면 Sharpe 큰 순으로 정렬해 첫 줄(\ft{iloc[0]})을 고른다.
\end{itemize}
\begin{workout}[따라 해 보기: \ft{snippets.py onese}]
\begin{lstlisting}[style=out]
best A 0.1: 3.8 - SE 0.8 = bar 3.0
pass: [['A', 0.1], ['A', 0.5], ['B', 0.5]]
pick: A 0.5
\end{lstlisting}
\Ans{최고는 A 10\%지만 기준선 3.0 이상인 후보 중 승인비율이 가장 큰 50\%에서, Sharpe가 더 높은 A를 고른다. 실제 결과(train)도 같은 구조였다. 최고는 부도확률 순 10\%(3.82, SE 0.84), 기준선은 2.98, 최종 제안은 부도확률 순 80\%(3.06)다.}
\end{workout}
\begin{whybox}[왜 최고값을 그냥 고르지 않나]
18개를 비교하면 운 좋게 높게 나온 후보가 1위가 되기 쉽다(다중 비교). SE가 0.7--0.8인데 1위와 2위의 차이가 0.2라면 그 순위는 의미가 없다. 구분되지 않는다면 많이 승인하는 쪽이 금액과 NPV가 크고, 모델이 극단 꼬리를 맞혀야 하는 부담도 적다. 머신러닝에서 규제 강도를 고를 때 쓰는 ``one-standard-error rule''과 같은 생각이다.
\end{whybox}

\section{파레토: 41번의 \texttt{dominated}}
\begin{lstlisting}
def dominated(mean, std):
    m, s = np.asarray(mean), np.asarray(std)
    return np.array([bool(((m >= m[i]) & (s <= s[i]) & ((m > m[i]) | (s < s[i]))).any())
                     for i in range(len(m))])
\end{lstlisting}
\textbf{3번 줄}은 후보 $i$마다 ``평균은 같거나 높고, 표준편차는 같거나 낮고, 둘 중 하나는 엄격히 나은 다른 후보가 하나라도(\ft{.any()}) 있나''를 묻는다. 있으면 $i$는 지배당한다(둘 다 나쁜 조합이라 뺀다).
\begin{workout}[따라 해 보기: \ft{snippets.py pareto}, 숫자는 train 승인 30\%]
\begin{lstlisting}[style=out]
ER frontier
lowPD dominated
simple frontier
random dominated
\end{lstlisting}
\Ans{단순 규칙(평균 1.84, 표준편차 0.48)이 부도확률 순(1.66, 0.50)보다 수익도 높고 흔들림도 작아 부도확률 순을 지배한다. 기대수익 순(1.88, 0.65)은 수익이 가장 높아 경계에 남는다. 다만 단순 규칙은 train을 보고 고른 규칙이라 test에서 확인해야 한다.}
\end{workout}

\section{test 실행: 40 \ensuremath{\to} 35 \ensuremath{\to} 15}
\begin{lstlisting}[language=bash,numbers=none]
python scripts/40_test_evaluate_v2_2.py --dry-run          # rehearsal on train
python scripts/40_test_evaluate_v2_2.py data/raw/test.csv  # the real test, once
\end{lstlisting}
40번은 35번 전체를 실행하고(35번은 다시 15번을 실행한다), 그 위에 v2.2 판정을 덧붙인다. 그래서 한 번 실행으로 v1, v2.1, v2.2가 모두 나온다. 결과는 \ft{results/test_v2_2/}에 생긴다.
\begin{table}[H]\centering\small
\begin{tabularx}{\linewidth}{@{}p{45mm}L@{}}\toprule
파일 & 무엇을 보나\\\midrule
\ft{summary.md} & 전체 요약. 맨 아래가 v2.2 판정\\
\ft{verdict.csv} & v1 가설 H1--H4 지지 여부\\
\ft{verdict_v2.csv} & v2.1 가설 H5--H9, 위약 점검\\
\ft{verdict_v2_2.csv} & H10(제안 유지), H11, H12, V3 도전자, H13(파레토)\\
\ft{V5_candidates_v2_2.csv} & test에서 18개 후보의 Sharpe와 순위\\
\ft{V6_dominance_v2_2.csv} & 같은 승인비율 지배 관계\\\bottomrule
\end{tabularx}
\end{table}
\begin{warnbox}[test는 한 번만]
결과를 보고 코드나 규칙을 고치면 test는 더 이상 증거가 아니다. 고쳐야 할 것이 생기면 \ft{docs/test_protocol_v2_2.md}의 ``이탈 기록''에 날짜, 내용, 이유를 적고, 고치기 전 결과도 함께 보고한다.
\end{warnbox}

\section{막힐 때}
\begin{table}[H]\centering\small
\begin{tabularx}{\linewidth}{@{}p{55mm}L@{}}\toprule
오류 메시지 & 뜻과 해결\\\midrule
\ft{command not found: python} & conda 환경이 꺼져 있다. \ft{conda activate cmapss}\\
\ft{ModuleNotFoundError: lc} & \ft{import _path}가 빠졌거나 \ft{scripts/} 밖에서 파일을 옮겼다\\
\ft{KeyError: 'xxx'} & 그런 열 이름이 없다. \ft{print(df.columns)}로 확인\\
\ft{ValueError: substring not found} & \ft{exec}로 자르는 기준 문자열이 앞 파일에서 바뀌었다\\
\ft{IndentationError} & 들여쓰기 칸 수가 섞였다(탭과 공백). 공백 4칸으로 통일\\
결과 숫자가 문서와 조금 다름 & seed, 승인 컷 단위(발행월인지 연도인지), 허들을 확인\\\bottomrule
\end{tabularx}
\end{table}

\end{document}
